
\section{Threats to validity}
\subsection{Internal validity}
The split is random stratified, not patient-stratified: cells from the
same patient can appear in train and test. The published reference uses
the same convention (its patient-level number, 95.9\%, is reported
separately in the source), so our cell-level comparison is like-for-like;
a patient-stratified re-split is the first robustness follow-up and is
flagged honestly here. Augmentation is seeded per sample; evaluation runs
exactly once per model on the test partition.
\subsection{Construct validity}
The census measures statistical label inconsistency under the
confident-learning model, not clinical truth. The admissibility gate
bounds when the estimator may run at all; the cleanlab cross-check bounds
implementation error (our flags are a strict subset of cleanlab's); the
forensics chapter bounds the alternative `bad acquisition' explanation
(only contrast separates the groups, small effect).
\subsection{External validity}
One collection, one stain, one scanner family. The BBBC battery and the
MedMNIST atlas bound how the method (not the model) transfers; the model
itself makes no cross-collection claim.
\subsection{Reliability}
Every number regenerates from committed probability dumps; two independent
metric libraries agree; the paper's tables are emitted programmatically
and a missing value renders UNVERIFIED rather than being filled.

\section{Census methodology in detail}
\subsection{Fold protocol}
Three folds, stratified, seeded; each fold's model trains from scratch on
the union of the other two (batch 64, Adam $10^{-3}$, early stopping,
patience 2). Out-of-fold probabilities cover every training cell exactly
once; the fold checkpoints are retained under
\texttt{results/malaria/oof\_ckpt} for audit.
\subsection{Threshold computation}
Per class $k$, the threshold $t_k$ is the mean self-confidence of
examples labelled $k$ (Definition~2.5); the confident joint counts
$(y_n, \arg\max \hat p)$ pairs with $\hat p_j \ge t_j$; off-diagonal mass
above the per-class expectation is the flag set. Our estimator and
cleanlab 2.x agree on the flag set direction (ours $\subset$ cleanlab's);
the implementation difference is threshold handling at ties.
\subsection{Why three folds}
Compute: one fold costs roughly one baseline training run; three is the
largest count that keeps the census inside the nightly budget. The
variance cost of fewer folds is acknowledged; the cleanlab cross-check
uses the same folds, so implementation comparison is fold-matched.
